\section{Sixteen unrestricted integers by a cofinal square radix}
\label{sec:signed16}

The new integer construction combines three components that were previously
separate: package 30's compact evaluation of the Jones ratio for its
31-variable quartic; package 29's signed relation combiner; and a new
restriction of the coding radix to a cofinal family of scaled squares.
The last step removes one positivity witness.  A smaller signed separator
also lowers the degree and evaluation cost.  The result is a completely
specified polynomial in \emph{sixteen unrestricted integers}, with
\[
 \deg P_{16}=813056\,5^{32}+10024,
 \qquad \operatorname{ops}(P_{16})=1865.
\]
These are improvements over the delivered 17-integer circuit, not a new
general bound for arbitrary recursively enumerable sets.  In particular,
Sun's 2026 ten-integer theorem is a different, more general representation
theorem, whose compiler is not instantiated here.

\subsection{Fixed index and literal Jones expressions}

Use the literal 31-variable quartic of package 30, supplied in the archive
as inert source data.  Its residuals are quadratic; shifting every positive
coordinate by one and summing their squares gives a normalized polynomial
\(P_{31}(x,\mathbf n)\), independent of the dummy input \(x\).
The source's fifth residual is \(EX+E-aX\), which equals one at the
all-ones tuple, so \(P_{31}(x,\mathbf0)>0\).  Its coefficient of the
fourth power of the retained first Pell root is one; hence its degree is
exactly four.  The quartic is rebound to the \emph{whole} coefficient tuple
proved in the coefficient section; its underlying four-port interface is
unchanged.  This is a semantic use of that interface theorem, not an
inference from matching file hashes.

Put \(\nu=31\) and \(L=5^{32}\).  If \(p_\alpha\) is the coefficient
of the exponent vector \(\alpha\in\mathbb N^{32}\), \(|\alpha|\le4\),
in the shifted quartic, define
\[
 c_\alpha=(4-|\alpha|)!\prod_{r=0}^{31}\alpha_r!\,p_\alpha,
 \qquad e(\alpha)=\sum_{r=0}^{31}\alpha_r5^r.
\]
There are 402 nonzero coefficient records, each a specified polynomial in
the four fixed compiler ports.  The bit lengths of those ports, in order
\((\beta,K_0,M_C,M_F)\), are
\[
 (1064921,1070004,1064921,1064922).
\]
For a coefficient record \(c_\alpha=\sum_t a_t\prod_j C_j^{e_{tj}}\),
the finite exact bound
\[
 |c_\alpha|\le H_\alpha
 :=\sum_t|a_t|2^{\sum_j e_{tj}\operatorname{bits}(C_j)},
 \qquad H_\alpha+1<2^{4269855}
\]
is verified for every record by the supplied index certificate.  Unlisted
records have coefficient zero.  Consequently the following are admissible
fixed integers:
\begin{align*}
 z_*&=2^{4269855}, &
 e_*(T)&=\sum_{|\alpha|\le4}(z_*+c_\alpha)
                          T^{L-e(\alpha)},\\
 y_*&=e_*(2z_*),&
 b_*&=(2z_*)^{L+1}=2^{4269856(L+1)}.
\end{align*}
The corresponding unused Jones index coordinate is
\(u_*=\sum_{r=1}^{31}(2z_*)^{5^r}\).  Base-five uniqueness makes the
exponents in \(e_*\) distinct, and every digit lies strictly between one
and \(2z_*\).  Thus
\[
 2(2z_*)^L\le y_*<(2z_*)^{L+1}=b_*.
\]
The integer \(y_*\) is needed to state the imported index theorem but
\emph{does not occur as a coefficient leaf of the new polynomial}.
It is replaced there by the simple power-of-two coefficient \(b_*\).
None of these fixed integers is an existential parameter.

The eight base variables are
\[
 (e,g_0,h_0,s_0,w_0,\phi,i_0,j_0)\in\mathbb Z^8.
\]
Set \(g=g_0+1\), \(h=h_0+1\), \(s=s_0+1\), \(w=w_0+1\),
\(i=i_0+1\), \(j=j_0+1\), and define
\begin{align*}
 b&=b_*e^2,& \mathcal B&=b^5,&q&=\mathcal B^L,\\
 \ell&=\sum_{r=1}^{31}\mathcal B^{5^r},&
 \lambda&=\sum_{r=0}^{4L-1}\mathcal B^r,&
 N&=q^{16},\\
 \theta&=\mathcal B-2z_*,&
 G&=32b\mathcal B^{5^{31}}.&
\end{align*}
Write \(E_*=e_*(\mathcal B)\), and put
\begin{align*}
 R={}&\bigl[g+E_*q^3+\ell q^5+
 \{2(E_*-z_*\lambda)(1+\mathcal B+g)^4
       +\lambda\mathcal B+\lambda\mathcal Bq^4\}q^7\bigr](N^2-N)\\
 &+\bigl[q^3-b\ell+\ell+\theta\lambda q^3+
              (\mathcal B-2)q^8\bigr](N^2-1).
\end{align*}
Continue with the literal Jones ratio expressions:
\begin{align*}
 U&=N^2w,&Y&=N^2s,&M&=RY,\\
 P&=2M^2U,&K&=R+1+h(P-1),&A&=M(U+1),\\
 B_1&=2R+1,&W_1&=bw,&C&=B_1+W_1+\phi,\\
 J&=4A-5,&Q&=3W_1C-2(W_1^2-1),&D&=(A^2-1)C^2+1,\\
 E&=iJDC^2,&F&=(A^2-1)E^2+1,&G_1&=A+F(CD+1-A),\\
 H&=B_1+jC,&I&=(G_1^2-1)H^2+1.&
\end{align*}
Finally set
\begin{align*}
 a_0&=(P^2-1)K^2+1,&a_1&=DFI,\\
 D_0&=JF,&T_0&=QF+J(H-C),\\
 m&=K^2G-4(C-KY)^2G-8gK^2.&
\end{align*}
Every intermediate symbol is a polynomial abbreviation.  In particular,
all exponents are fixed natural numbers, and no radical or variable
exponent occurs in the delivered arithmetic circuit.

\subsection{The new cofinal-radix lemma}

\paragraph{Lemma (restriction to scaled squares).}
For the fixed admissible Jones index, the source existential predicate
\[
 a_0,a_1\in\operatorname{Sq}(\mathbb Z),\qquad
 D_0\mid T_0,\qquad m>0
 \tag{JR}\label{eq:JR16}
\]
is equivalent to the parent-quartic zero predicate when
\(b=b_*e^2\), \(e\in\mathbb Z\setminus\{0\}\), and the other
seven source coordinates are nonnegative.

\emph{Proof.}
Soundness is immediate from \(b\ge b_*>y_*\): the integer
\(\varepsilon=b-y_*-1\) is nonnegative, and replacing the usual source
definition \(b=y_*+1+\varepsilon\) by this value changes no ratio
expression.  The source soundness theorem applies literally.

For completeness, inspect the constructive proof of
\texttt{Jones1982.USys\_of\_mem} in the pinned
\texttt{Master3.lean}.  For a fixed parent solution \(\mathbf n\),
its choice \(b=2^v\) uses the equality
\(v=x y_*+\sum_r n_r+1\) only to establish the initial bounds
\(v\ge1\), \(x y_*<b\), and \(n_r<b\).
After these bounds, the entire construction uses only those inequalities,
the positive power-of-two condition on \(b\), and the fixed index.
The displayed source was audited for every occurrence of that equality:
there is no later use.  Hence the same ordinary mathematical construction
works for every integer
\(v\ge x y_*+\sum_r n_r+1\).
Write \(b_*=2^\kappa\), where \(\kappa=4269856(L+1)\).
Choose \(t\ge0\) with \(\kappa+2t\) exceeding this bound, take
\(v=\kappa+2t\), and put \(e=2^t\).
Then \(b=2^v=b_*e^2\).  The unchanged
\texttt{Nine.reduced\_of\_mem} construction and
\texttt{Nine.ratio\_complete} give the ratio witnesses.
This is an explicit strengthening of the source construction; no
claim that an arbitrary Diophantine witness may be replaced by a square
is used. \hfill\(\square\)

\subsection{Five base guards and one margin guard}

Only \(\mathbf x=(g_0,s_0,w_0,\phi,j_0)\) needs direct
nonnegativity guards.  The variables \(h_0,i_0\) remain signed.
Introduce six integers \(z_1,\ldots,z_6\).  Put \(C_0=1\) and
\begin{align*}
 d_r&=3x_r+2,& B_r&=1+4d_r C_{r-1}^2z_r^2,&
 C_r&=C_{r-1}B_r\quad(1\le r\le5),\\
 d_6&=3m-1,&B_6&=1+4d_6 C_5^2z_6^2,&C_*&=C_5B_6.
\end{align*}

\paragraph{Lemma (coprime Pell guard).}
If every \(z_r\ne0\), then
\[
 C_*\in\operatorname{Sq}(\mathbb Z)
 \quad\Longrightarrow\quad
 g_0,s_0,w_0,\phi,j_0\ge0,\quad m>0.
\]
Conversely these inequalities permit a choice of nonzero \(z_r\)
for which \(C_*\) is a square.  Moreover, for arbitrary signed
base coordinates and nonzero \(z_r\),
\(C_r\ne0\) and \(|C_*|\ge3^6\).

\emph{Proof.}
Each \(d_r\equiv2\pmod3\) is nonzero, and each
\(B_r\equiv1\pmod4\) is nonzero.  Also
\(\gcd(B_r,C_{r-1})=1\), so the factors \(B_r\) are pairwise
coprime.  If their product is a square, unique factorization makes each
\(|B_r|\) a square.  A negative square is zero or three modulo four;
therefore every \(B_r\) is a positive square.
A negative \(d_r\), with both \(C_{r-1},z_r\ne0\), would give
\(B_r\le-3\).  Thus \(d_r>0\), giving the asserted inequalities.

Conversely, each positive \(d_r\equiv2\pmod3\) is nonsquare.
The Pell equation \(u^2-d_r v^2=1\) has a nontrivial positive solution
whose \(v\) is divisible by any specified positive integer: reduce a
Pell unit matrix modulo that integer and take a positive power equal
to the identity in the finite matrix group.  Take the modulus
\(2|C_{r-1}|\) and set \(z_r=v/(2C_{r-1})\).
This makes \(B_r=u^2\) and \(z_r\ne0\).  Induction proves
completeness.  Finally, without any positivity assumption,
\(|B_r|\ge3\) when \(z_r\ne0\); multiply the six inequalities.
\hfill\(\square\)

\paragraph{Lemma (the two unguarded base coordinates).}
Assume \(e\ne0\), the five base guards are nonnegative,
\(a_0\) is a square, and \(m>0\).  Then \(h\ge1\).
If also \(i\ne0\), replacing \(i\) by \(|i|\) preserves
\eqref{eq:JR16} and makes \(i_0=|i|-1\) nonnegative.

\emph{Proof.}
Here \(b\ge b_*>y_*\), so \(\mathcal B>2z_*\),
\(G>0\), and \(g\ge1\).  From \(m>0\) obtain
\[
 K\ne0,\qquad g<G/8=4b\mathcal B^{5^{31}},\qquad
 4(C-KY)^2<K^2.
\]
Write \(r=5^{31}\).  The first two bounds imply
\(1+\mathcal B+g<\mathcal B^{r+2}\le q\).
Also \(\ell<\mathcal B^{r+1}\), so \(b\ell<q^3\).
In the first bracket defining \(R\), the positive term
\(\lambda\mathcal Bq^4\) strictly absorbs
\(2z_*\lambda(1+\mathcal B+g)^4\); every remaining term is
nonnegative.  The second bracket is positive as well.  Thus
\(R\ge N^2-N\ge2\).
Now \(Y,U\ge1\), \(C=2R+1+bw+\phi>0\), and
\(P=2R^2Y^2U>R+2\).
If \(K<0\), then \(|C-KY|>|K|\), contradicting the displayed
strict inequality.  Hence \(K>0\).
If \(h\le0\), then \(1\le K\le R+1<P\).
The case \(K=1\) would imply \(P-1\mid R\), impossible.
For \(2\le K<P\),
\[
 (PK-1)^2<(P^2-1)K^2+1<(PK)^2,
\]
contradicting squarehood of \(a_0\).  This proves \(h\ge1\).
Finally \(i\) occurs only linearly in \(E\), whose only relevant
consumer is \(E^2\) in \(F\).  Changing the sign of \(i\)
therefore changes none of \(a_0,a_1,D_0,T_0,m\).
\hfill\(\square\)

\subsection{Divisibility and nonvanishing with one extra coordinate}

Use the fixed intersective polynomial
\[
 \Pi(t)=(t^3-t-1)(t^2+23).
\]
It has no rational root.  It has a root modulo every positive integer:
for odd primes other than 23, either the cubic has a simple root, or
it is irreducible and its Frobenius three-cycle makes its discriminant
\(-23\) a nonzero square, giving a simple root of the quadratic.
These roots lift to all prime powers.  At 23, the cubic has the simple
root 3; at 2, \(-23\equiv1\pmod8\) has square roots modulo every
power of two.  The Chinese remainder theorem finishes the proof.

Define the homogeneous degree-five polynomial
\[
 \mathcal R(D,T,v)=
 ((T+Dv)^3-D^2(T+Dv)-D^3)((T+Dv)^2+23D^2).
\]
For integers \(D\ne0,T,Z\),
\[
 (D\mid T\ \text{and}\ Z\ne0)
 \quad\Longleftrightarrow\quad
 \exists v\in\mathbb Z:\ ZD^5\mid\mathcal R(D,T,v).
 \tag{IM}\label{eq:IM16}
\]
Indeed \(\mathcal R=D^5\Pi(T/D+v)\) never vanishes, since
\(T/D+v\) is rational.  If \(D\mid T\), intersectivity supplies
\(v\) so that \(Z\mid\Pi(T/D+v)\).  Conversely, if a prime
satisfies \(\alpha=v_p(D)>\beta=v_p(T)\), then
\(v_p(T+Dv)=\beta\), and the two factors in \(\mathcal R\)
have unique lowest-valuation terms of valuations \(3\beta\)
and \(2\beta\).  Hence
\(v_p(\mathcal R)=5\beta<5\alpha\), contradicting divisibility
by \(D^5\).  The case \(Z=0\) is excluded by nonvanishing.

For the Jones expressions, \(D_0=JF\ne0\) at \emph{every}
integer assignment: \(J=4A-5\ne0\); and if \(F=0\), then
\((1-A^2)E^2=1\) gives \(A=0,E=\pm1\), although
\(E=iJDC^2\) would then be divisible by five.  This contradiction
does not use any guard or desired square condition.

\subsection{The literal sixteen-variable polynomial}

Besides the eight base coordinates and six \(z_r\), introduce
\(v,u\in\mathbb Z\).  These are all sixteen unknowns.  Put
\begin{align*}
 Z&=ie\prod_{r=1}^{6}z_r,& S&=ZD_0^5,&T&=\mathcal R(D_0,T_0,v),\\
 w_1&=C_*,&w_2&=C_*(P^2+K^2+2),&X&=Su-T,\\
 A_0&=S^2C_*,& A_1&=S^2w_1^2a_0,& A_2&=S^2w_2^2a_1,\\
 V&=X^2+A_2-A_0-A_1,&W&=V^2+4X^2A_2-4A_0A_1.&
\end{align*}
The equation is
\[
 \boxed{\quad P_{16}=W^2-16V^2X^2A_2=0.\quad}
\]

\paragraph{Lemma (signed three-root finalizer).}
A zero of this polynomial gives
\[
 S\ne0,\qquad S\mid T,\qquad
 C_*,a_0,a_1\in\operatorname{Sq}(\mathbb Z).
 \tag{F}\label{eq:F16}
\]
Conversely, these conditions on the earlier coordinates allow an
integer choice of \(u\) making the polynomial zero.

\emph{Proof.}
If \(S=0\), then \(A_0=A_1=A_2=0\), \(X=-T\), and
\(P_{16}=T^8\ne0\).  Thus any zero has \(S\ne0\), so
\(i,e,z_1,\ldots,z_6\) are nonzero.  The guard lemma, before
any squarehood conclusion, gives \(|C_*|\ge3^6\).
Consequently
\[
 |w_1|>\sqrt{|C_*|}.
\]
For arbitrary real \(P,K\),
\[
 \sqrt{|(P^2-1)K^2+1|}\le P^2+K^2+1.
\]
To verify this, first bound the absolute radicand by
\((P^2+1)K^2+1\), then subtract this expression from
\((P^2+K^2+1)^2\); the result is
\(P^4+K^4+P^2K^2+2P^2+K^2\ge0\).
It follows that
\[
 |w_2|>\sqrt{|C_*|}+|w_1|\sqrt{|a_0|}.
\]
These estimates permit signed \(C_*,w_1,w_2\); positivity has
not been assumed.

Choose complex square roots \(r_0^2=C_*\), \(r_1^2=a_0\),
\(r_2^2=a_1\).  Direct multiplication gives the polynomial identity
\[
 P_{16}=\prod_{\epsilon_0,\epsilon_1,\epsilon_2\in\{-1,1\}}
 \bigl(X-S(\epsilon_0r_0+\epsilon_1w_1r_1+
                                \epsilon_2w_2r_2)\bigr).
\]
At a zero, one of these signed sums is rational.  Under any
automorphism of the finite Galois extension generated by the three
roots, each root changes only by sign.  If the largest nonzero root
whose sign changes has index 1 or 2, its weighted modulus strictly
exceeds the sum of all earlier weighted moduli by the preceding
bounds: a nonzero root of an integer radicand has modulus at least
one.  A change only at index 0 is impossible as well.  Thus an
automorphism cannot preserve the rational sum while changing a
nonzero root.  Every root is fixed, hence rational, and a rational
square root of an integer is an integer.  This handles negative
candidate radicands, zero roots, and dependent square classes.
The resulting integral signed sum equals \(u-T/S\), so
\(S\mid T\).  This proves necessity.
Conversely, select integral roots and take
\(u=T/S+r_0+w_1r_1+w_2r_2\).  A factor vanishes.
\hfill\(\square\)

\paragraph{Theorem (the concrete integer equation).}
With the fixed coefficients above,
\[
 \mathrm{ZFC}\vdash
 \bigl(\nexists\mathbf t\in\mathbb Z^{16}:
                           P_{16}(\mathbf t)=0\bigr)
 \longleftrightarrow\operatorname{Con}(\mathrm{ZFC}).
\]

\emph{Proof.}
A zero first gives \eqref{eq:F16}.  Apply \eqref{eq:IM16} to
recover \(D_0\mid T_0\) and all nonzero conditions.
The coprime guard gives the five nonnegative base coordinates and
\(m>0\).  The unguarded-coordinate lemma gives \(h_0\ge0\)
and permits replacing \(i\) by \(|i|\).
Because \(e\ne0\), the derived source coordinate
\(\varepsilon=b_*e^2-y_*-1\) is nonnegative.  We have therefore
recovered the literal source Jones predicate, which gives a zero of
the parent quartic and hence, through the whole-tuple compiler and
the ordinary K215 calibration, a finite ZFC refutation.

Conversely, a finite refutation gives a parent-quartic zero.
The cofinal-radix lemma gives Jones witnesses with
\(b=b_*e^2\), \(e\ne0\).  The Pell guard supplies the six
nonzero \(z_r\) and a square \(C_*\).
Intersectivity supplies \(v\) with \(S\mid T\), and the
finalizer supplies \(u\).  All steps are ordinary number theory
and finite Galois theory formalizable in ZFC.  Negating this
two-way existential equivalence proves the assertion.
It proves neither consistency nor inconsistency. \hfill\(\square\)

\subsection{Exact complexity and reproducible checks}

The square-radix substitution doubles the degree contributed by the
radix coordinate, while removing one guard.  With \(L=5^{32}\),
selected exact degrees are
\[
\begin{array}{c|r@{\qquad}c|r}
 R&470L&P&1900L+3\\
 K&1900L+4&m&4442L+12\\
 a_0&7600L+14&a_1&49280L+68\\
 D_0&13750L+20&C_*&4442L+1103\\
 S&68750L+108&T&68750L+105\\
 X&68750L+109&w_2&8242L+1111\\
 A_2&203264L+2506&P_{16}&813056L+10024
\end{array}
\]
The degree of \(A_2\) strictly exceeds those of
\(X^2,A_0,A_1\).  Therefore \(V\) has leading form \(A_2\),
\(W\) has leading form \(A_2^2\), and the leading form of the
final output is \(A_2^4\).  On the diagonal specialization this is
a positive integer times a power of the positive fixed coefficient
\(b_*\).  The supplied leading-term certificate keeps that power
symbolic, establishing the stated exact degree for the authentic
tuple rather than only a generic upper bound.

All fixed powers are expanded into binary multiplications; squaring
and scalar multiplication are charged.  Reuse is free, and fixed
coefficient construction is excluded from variable-evaluation cost.
The emitted, fully live DAG has 705 additions, 26 subtractions,
and 1134 multiplications, totaling 1865.  Its source digit polynomial
uses the compact Horner evaluator and shared radix-five powers of
package 30; no imported Python program is executed.

The independent diagnostics evaluate actual fixed coefficients modulo
three primes and compare 90 signed tuples both with a direct formula
and with the unchanged inherited Jones ratio after substituting
\(\varepsilon=b_*e^2-y_*-1\).  They also check 200 independent
eight-factor norm products on integral roots and another 200 products
in the rank-eight free radical algebra with arbitrary signed radicands,
the basic separator inequalities on
57,222 assignments, and intersectivity modulo every integer through
1200.  These finite tests check definitions and arithmetic.  The
unbounded correctness argument is the proof above, together with
the explicitly identified source compiler and Jones theorems; the
archive does not claim a new end-to-end Lean or Rocq certificate.

\paragraph{Seventeen-integer companion.}
Retaining \(b=y_*+1+\varepsilon\), guarding all six base
coordinates as in package 29, and using the improved separator
\(w_1=C_*\) gives a companion polynomial in seventeen integers.
Its exact degree is \(406528\,5^{32}+27512\), and its 1874 live
operations consist of 709 additions, 26 subtractions, and 1139
multiplications.  It uses the same new coefficient tuple and index.
The corresponding proof is the signed-three-root proof above with
seven guard factors and without the cofinal-square-radix step.
This is an explicit alternative when degree is valued more highly
than the final integer-variable reduction.
